\clearpage
\section{Introduction}
\label{sec:Intro}

Higher education is widely regarded as an important instrument for promoting intergenerational mobility and reducing earnings inequality. Yet whether expansions in university access translate into broad and equitable labour-market gains remains an open empirical question, particularly in developing economies marked by strong regional and demographic disparities. Brazil---famously described as \emph{Belíndia}, a country that simultaneously resembles Belgium and India \cite{bacha1976}---offers a particularly compelling setting: its labour market exhibits persistent wage gaps by ethnicity, gender, and region even after controlling for observable human capital \cite{AriasYamadaTejerina2002, Telles2004, deoliveira2021}.

In 2005, the Brazilian federal government launched the \acrfull{prouni}, a large-scale scholarship programme that exchanges tax exemptions for private \acrfull{hei} participation in the provision of tuition-free university places for low-income students. By 2019, ProUni had allocated more than 3.5 million scholarships, and its institutional design included affirmative-action provisions reserving places for self-declared non-white applicants in proportion to their state-level population share \cite{MEC2024}. The programme therefore constitutes both an expansion in higher-education access and a redistributive intervention aimed at historically disadvantaged groups.

Despite its scale and policy importance, the labour-market effects of ProUni remain insufficiently understood. Existing work has documented positive effects of affirmative action in elite public universities \cite{RiehlMachadoReyes2023} and estimated returns to education in Brazil using more conventional Mincerian approaches \cite{Card1999, Mincer1974}, but there is still limited causal evidence on how ProUni's geographically uneven expansion affected local labour markets. In particular, little is known about whether the programme generated similar gains across demographic groups and across regions with very different higher-education capacity and labour-demand conditions.

This thesis addresses that gap by asking: \emph{How does the expansion of ProUni scholarships affect formal-sector earnings, and do these effects vary systematically by gender, ethnicity, and geographic region?}

To answer this question, I construct a panel dataset linking three administrative sources---ProUni scholarship records, the \acrfull{rais} employer--employee census, and Census population data---at the microregion level over the period 2005--2019. The empirical strategy builds a lagged cumulative scholarship-density measure intended to proxy the local exposure of labour markets to ProUni-supported graduates, normalized by the youth population. Because the effects of educational expansion may be non-linear, I discretize this continuous exposure measure into dose strata and estimate heterogeneous treatment effects separately within those strata.

Identification relies on the heterogeneity-robust \acrfull{did} estimator of \citeA{CallawaySantAnna2021}, which is designed for settings with staggered treatment timing and potentially heterogeneous effects, and therefore avoids the well-known weighting problems of conventional \acrfull{twfe} regression \cite{GoodmanBacon2021, deChaisemartin2020}. The event-study implementation makes it possible to inspect pre-treatment dynamics transparently and to distinguish the primary causal objects of the analysis---within-stratum group-time average treatment effects---from more tentative cross-stratum comparisons. For the latter, the thesis follows the conceptual guidance of \citeA{CallawayGoodmanBaconSantAnna2024} and interprets dose contrasts cautiously.

The thesis makes three contributions. First, it provides new reduced-form evidence on the local labour-market consequences of ProUni using a research design that exploits both temporal and spatial variation in programme intensity. Second, by estimating separate models by gender, ethnicity, and broad region, it examines the distributional profile of a large-scale access policy rather than treating the programme's effects as homogeneous. Third, it contributes an applied case study of how discretized treatment-intensity designs can be used in modern difference-in-differences frameworks when the underlying policy exposure is continuous but highly skewed.

The remainder of the thesis is organized as follows. Section~\ref{sec:literature} reviews the theoretical foundations in human capital and segmented labour-market theory, together with the empirical literature on returns to education in Brazil and internationally. Section~\ref{sec:institutional_context} presents the institutional context, covering Brazil's higher-education system, the design and legal foundations of ProUni, and concurrent policy interventions. Section~\ref{sec:methodology} details the data sources, treatment construction, identification assumptions, and estimation strategy. Section~\ref{sec:results} presents the main results. Section~\ref{sec:discussion} interprets the findings and their limitations. Section~\ref{sec:conclusion} concludes.
